\documentclass[11pt]{article}
\usepackage[utf8]{inputenc}
\usepackage[T1]{fontenc}
\usepackage{amsmath}
\usepackage{amssymb}
\usepackage{multicol}
\usepackage{geometry}
\usepackage{tikz}
\usepackage{enumitem}
\usepackage{xcolor}
\usepackage{titlesec}

% Layout
\geometry{a4paper, left=1cm, right=1cm, top=0.5cm, bottom=1.2cm}
\setlength{\columnseprule}{0.4pt}

% Títulos
\titleformat{\section}{\normalfont\Large\bfseries\color{blue}}{}{0em}{}
\titleformat{\subsection}{\normalfont\large\bfseries\color{red}}{}{0em}{}

\title{\textcolor{red}{Segundo Trimestre: Matemática
\\Razão e Proporção}}
\author{Professor: Jefferson}
\date{}

\begin{document}

\maketitle
\vspace{-1cm}

\begin{multicols}{2}

\section*{Questões Comentadas}

\subsection*{Questão 1}
Calcule a razão entre 18 e 6.

\textbf{Comentário:} Razão é uma divisão.
\[
\frac{18}{6} = 3
\]
\textbf{Resposta:} 3

\subsection*{Questão 2}
A razão entre dois números é $4:5$. Escreva na forma de fração.

\textbf{Comentário:} Razão $a:b = \frac{a}{b}$.
\[
\frac{4}{5}
\]
\textbf{Resposta:} $\frac{4}{5}$

\subsection*{Questão 3}
Verifique se $\frac{6}{8}$ e $\frac{9}{12}$ formam uma proporção.

\textbf{Comentário:}
\[
6 \cdot 12 = 72 \quad \text{e} \quad 8 \cdot 9 = 72
\]
\textbf{Resposta:} Sim, formam proporção.

\subsection*{Questão 4}
Determine $x$:
\[
\frac{3}{4} = \frac{x}{20}
\]

\textbf{Comentário:}
\[
3 \cdot 20 = 4x \Rightarrow x = 15
\]

\textbf{Resposta:} $x = 15$

\subsection*{Questão 5}
Se 7 cadernos custam R\$ 35, quanto custam 10 cadernos?

\textbf{Comentário:} Grandezas diretamente proporcionais.
\[
\frac{7}{10} = \frac{35}{x}
\Rightarrow x = 50
\]

\textbf{Resposta:} R\$ 50

\subsection*{Questão 6}
Um carro percorre 180 km em 3 horas. Quantos km percorrerá em 5 horas?

\textbf{Comentário:} Distância e tempo são diretamente proporcionais.
\[
\frac{180}{3} = \frac{x}{5} \Rightarrow x = 300
\]

\textbf{Resposta:} 300 km

\subsection*{Questão 7}
4 operários fazem uma obra em 15 dias. Quantos dias levarão 5 operários?

\textbf{Comentário:} Grandezas inversamente proporcionais.
\[
4 \cdot 15 = 5 \cdot x \Rightarrow x = 12
\]

\textbf{Resposta:} 12 dias

\subsection*{Questão 8}
Uma torneira enche um tanque em 6 horas. Quanto tempo levarão 3 torneiras?

\textbf{Comentário:}
\[
1 \cdot 6 = 3 \cdot x \Rightarrow x = 2
\]

\textbf{Resposta:} 2 horas

\subsection*{Questão 9}
Em uma receita, 4 ovos produzem 24 bolos. Quantos bolos serão feitos com 6 ovos?

\textbf{Comentário:} Grandezas diretamente proporcionais.
\[
\frac{4}{6} = \frac{24}{x} \Rightarrow x = 36
\]

\textbf{Resposta:} 36 bolos

\subsection*{Questão 10}
Um mapa está na escala 1:50000. Qual a distância real de 4 cm no mapa?

\textbf{Comentário:}
\[
4 \cdot 50000 = 200000 \text{ cm} = 2 \text{ km}
\]

\textbf{Resposta:} 2 km

\subsection*{Questão 11}
Se 12 lápis custam R\$ 18, quanto custam 20 lápis?

\textbf{Comentário:}
\[
\frac{12}{20} = \frac{18}{x} \Rightarrow x = 30
\]

\textbf{Resposta:} R\$ 30

\subsection*{Questão 12}
6 máquinas produzem uma peça em 10 horas. Em quanto tempo 3 máquinas farão o mesmo trabalho?

\textbf{Comentário:} Inversamente proporcionais.
\[
6 \cdot 10 = 3 \cdot x \Rightarrow x = 20
\]

\textbf{Resposta:} 20 horas

\subsection*{Questão 13}
A razão entre o número de meninas e meninos é $5:4$. Quantos meninos há se existem 20 meninas?

\textbf{Comentário:}
\[
\frac{5}{4} = \frac{20}{x} \Rightarrow x = 16
\]

\textbf{Resposta:} 16 meninos

\subsection*{Questão 14}
Determine $x$:
\[
\frac{x}{9} = \frac{14}{21}
\]

\textbf{Comentário:}
\[
21x = 126 \Rightarrow x = 6
\]

\textbf{Resposta:} $x = 6$

\subsection*{Questão 15}
Quanto maior a velocidade, menor o tempo para percorrer uma distância fixa. Essa relação é:

\textbf{Comentário:} Uma aumenta e a outra diminui.

\textbf{Resposta:} Inversamente proporcional

\subsection*{Questão 16}
5 metros de tecido custam R\$ 40. Quanto custam 8 metros?

\textbf{Comentário:}
\[
\frac{5}{8} = \frac{40}{x} \Rightarrow x = 64
\]

\textbf{Resposta:} R\$ 64

\subsection*{Questão 17}
2 pedreiros constroem um muro em 18 dias. Quantos dias 6 pedreiros levarão?

\textbf{Comentário:}
\[
2 \cdot 18 = 6 \cdot x \Rightarrow x = 6
\]

\textbf{Resposta:} 6 dias

\subsection*{Questão 18}
A razão entre 45 e 15 é:

\textbf{Comentário:}
\[
\frac{45}{15} = 3
\]

\textbf{Resposta:} 3

\subsection*{Questão 19}
Se 3 kg de açúcar custam R\$ 12, qual o preço de 1 kg?

\textbf{Comentário:}
\[
\frac{3}{1} = \frac{12}{x} \Rightarrow x = 4
\]

\textbf{Resposta:} R\$ 4

\subsection*{Questão 20}
É possível resolver problemas de razão e proporção usando regra de três?

\textbf{Comentário:} A regra de três é uma aplicação direta das proporções.

\textbf{Resposta:} Sim.

\end{multicols}

\end{document}

